\chapter{Conclusions and Future Work}\label{ch:conclusions}

\section{Principal conclusions}
This thesis develops a staged framework for physics-informed
fluid--particle transport. The progression begins with a one-way coupled
particle problem, introduces Gaussian parcel reaction in a straight
channel, and defines a smooth constriction for the next study. The
completed evidence supports the following conclusions.

First, the baseline carrier network agrees with an independent
finite-element solution of the stated curved-channel problem, with velocity
and gauge-aligned pressure errors of \FlowVelocity\% and
\FlowPressure\%. Analytical plane-flow verification and reference-mesh
sensitivity support this comparison. The normalized learned Stokes
response has an independent error of \DragError\%.

Second, carrier-field agreement does not by itself establish accurate
particle transport. The seven baseline comparisons do not all meet the
endpoint, position, and time-coverage conditions. Deposition and
wall-sensitive parameter trends are therefore outside the accepted
quantitative conclusions.

Third, a normalized Gaussian used for both interpolation and spreading
preserves discrete force and work on the transfer quadrature. The resulting
interphase drag power is nonpositive. These identities are algebraic
properties of the discretization and do not imply that an imperfectly
trained field satisfies every global balance.

Fourth, the reduced Gaussian-model tests recover Stokes response, uniform
two-way relaxation, kernel integrals, and Taylor--Green decay with the
reported errors and residuals. The accepted solution-observable errors
in the listed relaxation and vortex tests are below 0.25\%. The fully
coupled straight-channel predictions remain diagnostic because their
complete residual criteria are not met.

Finally, the constricted-channel extension has a consistent geometric and
evaluation definition, including physical-coordinate derivatives,
fluid-domain kernel normalization, and independent throat-resolution
checks. It has no completed predictive result in the present study.

\section{Contributions}
The study contributes a common evaluation framework connecting carrier
fields, drag approximation, transfer consistency, and particle outcomes.
It gives a dimensionally consistent planar sphere-parcel formulation and
explicitly separates a Gaussian force distribution from a resolved solid
body. It also defines which verified components can be retained when
extending the model to a variable-height channel.

These contributions provide a practical basis for further computation
without assigning unsupported accuracy to complex coupled predictions.
The numerical record and portable manuscript sources preserve the arrays,
criteria, and figure provenance needed to reassess the stated findings.

\section{Priorities for future work}
The immediate priority is to improve and independently verify the
transport-sensitive parts of the calculation. For the baseline, this
means near-wall field assessment and a direct exact-drag comparison. For
Gaussian coupling, it means a matched independent regularized channel
solution, transfer-quadrature refinement, temporal sampling refinement,
and multiple training seeds at fixed physical parameters. Sampling and
numerical-uncertainty methodologies provide a basis for planning those
comparisons \citep{wu2023,celik2008,oberkampf2002}. % Citation-plan C35

The constricted study should then begin with an unladen flow and a
straight-channel boundary-condition control. One-way particles can be
introduced after the carrier reference passes, followed by low-loading
Gaussian feedback. Physical deposition or dynamic clogging requires a
separately specified wall and population model. Experimental comparison
will require matched geometry, operating conditions, particle properties,
and observables before it can establish physical validity.
